% Options for packages loaded elsewhere
% Options for packages loaded elsewhere
\PassOptionsToPackage{unicode}{hyperref}
\PassOptionsToPackage{hyphens}{url}
\PassOptionsToPackage{dvipsnames,svgnames,x11names}{xcolor}
%
\documentclass[
  11pt,
]{article}
\usepackage{xcolor}
\usepackage[margin=1.0in]{geometry}
\usepackage{amsmath,amssymb}
\setcounter{secnumdepth}{-\maxdimen} % remove section numbering
\usepackage{iftex}
\ifPDFTeX
  \usepackage[T1]{fontenc}
  \usepackage[utf8]{inputenc}
  \usepackage{textcomp} % provide euro and other symbols
\else % if luatex or xetex
  \usepackage{unicode-math} % this also loads fontspec
  \defaultfontfeatures{Scale=MatchLowercase}
  \defaultfontfeatures[\rmfamily]{Ligatures=TeX,Scale=1}
\fi
\usepackage{lmodern}
\ifPDFTeX\else
  % xetex/luatex font selection
\fi
% Use upquote if available, for straight quotes in verbatim environments
\IfFileExists{upquote.sty}{\usepackage{upquote}}{}
\IfFileExists{microtype.sty}{% use microtype if available
  \usepackage[]{microtype}
  \UseMicrotypeSet[protrusion]{basicmath} % disable protrusion for tt fonts
}{}
\usepackage{setspace}
\makeatletter
\@ifundefined{KOMAClassName}{% if non-KOMA class
  \IfFileExists{parskip.sty}{%
    \usepackage{parskip}
  }{% else
    \setlength{\parindent}{0pt}
    \setlength{\parskip}{6pt plus 2pt minus 1pt}}
}{% if KOMA class
  \KOMAoptions{parskip=half}}
\makeatother
% Make \paragraph and \subparagraph free-standing
\makeatletter
\ifx\paragraph\undefined\else
  \let\oldparagraph\paragraph
  \renewcommand{\paragraph}{
    \@ifstar
      \xxxParagraphStar
      \xxxParagraphNoStar
  }
  \newcommand{\xxxParagraphStar}[1]{\oldparagraph*{#1}\mbox{}}
  \newcommand{\xxxParagraphNoStar}[1]{\oldparagraph{#1}\mbox{}}
\fi
\ifx\subparagraph\undefined\else
  \let\oldsubparagraph\subparagraph
  \renewcommand{\subparagraph}{
    \@ifstar
      \xxxSubParagraphStar
      \xxxSubParagraphNoStar
  }
  \newcommand{\xxxSubParagraphStar}[1]{\oldsubparagraph*{#1}\mbox{}}
  \newcommand{\xxxSubParagraphNoStar}[1]{\oldsubparagraph{#1}\mbox{}}
\fi
\makeatother


\usepackage{longtable,booktabs,array}
\usepackage{calc} % for calculating minipage widths
% Correct order of tables after \paragraph or \subparagraph
\usepackage{etoolbox}
\makeatletter
\patchcmd\longtable{\par}{\if@noskipsec\mbox{}\fi\par}{}{}
\makeatother
% Allow footnotes in longtable head/foot
\IfFileExists{footnotehyper.sty}{\usepackage{footnotehyper}}{\usepackage{footnote}}
\makesavenoteenv{longtable}
\usepackage{graphicx}
\makeatletter
\newsavebox\pandoc@box
\newcommand*\pandocbounded[1]{% scales image to fit in text height/width
  \sbox\pandoc@box{#1}%
  \Gscale@div\@tempa{\textheight}{\dimexpr\ht\pandoc@box+\dp\pandoc@box\relax}%
  \Gscale@div\@tempb{\linewidth}{\wd\pandoc@box}%
  \ifdim\@tempb\p@<\@tempa\p@\let\@tempa\@tempb\fi% select the smaller of both
  \ifdim\@tempa\p@<\p@\scalebox{\@tempa}{\usebox\pandoc@box}%
  \else\usebox{\pandoc@box}%
  \fi%
}
% Set default figure placement to htbp
\def\fps@figure{htbp}
\makeatother





\setlength{\emergencystretch}{3em} % prevent overfull lines

\providecommand{\tightlist}{%
  \setlength{\itemsep}{0pt}\setlength{\parskip}{0pt}}



 


\usepackage[left]{lineno}
\linenumbers
\modulolinenumbers
\usepackage{helvet}
\renewcommand*\familydefault{\sfdefault}
\usepackage[T1]{fontenc}
\makeatletter
\@ifpackageloaded{caption}{}{\usepackage{caption}}
\AtBeginDocument{%
\ifdefined\contentsname
  \renewcommand*\contentsname{Table of contents}
\else
  \newcommand\contentsname{Table of contents}
\fi
\ifdefined\listfigurename
  \renewcommand*\listfigurename{List of Figures}
\else
  \newcommand\listfigurename{List of Figures}
\fi
\ifdefined\listtablename
  \renewcommand*\listtablename{List of Tables}
\else
  \newcommand\listtablename{List of Tables}
\fi
\ifdefined\figurename
  \renewcommand*\figurename{Figure}
\else
  \newcommand\figurename{Figure}
\fi
\ifdefined\tablename
  \renewcommand*\tablename{Table}
\else
  \newcommand\tablename{Table}
\fi
}
\@ifpackageloaded{float}{}{\usepackage{float}}
\floatstyle{ruled}
\@ifundefined{c@chapter}{\newfloat{codelisting}{h}{lop}}{\newfloat{codelisting}{h}{lop}[chapter]}
\floatname{codelisting}{Listing}
\newcommand*\listoflistings{\listof{codelisting}{List of Listings}}
\makeatother
\makeatletter
\makeatother
\makeatletter
\@ifpackageloaded{caption}{}{\usepackage{caption}}
\@ifpackageloaded{subcaption}{}{\usepackage{subcaption}}
\makeatother
\usepackage{bookmark}
\IfFileExists{xurl.sty}{\usepackage{xurl}}{} % add URL line breaks if available
\urlstyle{same}
\hypersetup{
  colorlinks=true,
  linkcolor={blue},
  filecolor={Maroon},
  citecolor={Blue},
  urlcolor={Blue},
  pdfcreator={LaTeX via pandoc}}


\author{}
\date{}
\begin{document}


\setstretch{1.75}
\raggedright

\section{strollur: An R package for working with amplicon sequence data
in
R}\label{strollur-an-r-package-for-working-with-amplicon-sequence-data-in-r}

\vspace{20mm}

\textbf{Running title:} mums2

\vspace{20mm}

Sarah L. Westcott\textsuperscript{1}, Gregory Johnson
Jr.\textsuperscript{1}, Patrick D. Schloss\textsuperscript{1\textdagger}

\vspace{25mm}

\textdagger To whom correspondence should be addressed\\
\href{mailto:pschloss@umich.edu}{pschloss@umich.edu}

\vspace{10mm}

1 Department of Microbiology \& Immunology, University of Michigan, Ann
Arbor, MI 48109

\vspace{20mm}

\textbf{Software Announcement}

\newpage

\subsection{Abstract}\label{abstract}

Microbiologists are increasingly relying on R to analyze and visualize
amplicon sequence data from microbiome studies. We present strollur,
which isan open source package that facilitates the import, export, and
handling of amplicon sequence data within the R programming language.

\newpage

\subsection{Announcement}\label{announcement}

Amplicon sequencing projects including 16S rRNA gene sequence
collections are widely used by microbiologists to analyze diverse
microbiomes. These datasets are often large and complicated making it
difficult for researchers to keep their data organized and record the
provenance of derivative data products. A long-standing difficulty for
many has been the ability to move data between different analytical
tools. This often requires reformatting data, which is challenging for
researchers with limited programming skills. Several efforts have been
made to facilitate the management of amplicon sequencing data and their
associated metadata including the biom file format
{[}\textbf{\emph{ref}}{]}, mothur output files
{[}\textbf{\emph{ref}}{]}, the qiime2 archive {[}\textbf{\emph{ref}}{]},
and phyloseq R objects {[}\textbf{\emph{ref}}{]}. Each of these formats
and tools have strengths but also limitations that make their use
difficult for researchers new to microbiome data analysis or who are
more accustomed to using R's tidyverse collection of packages.

To overcome these challenges, we developed the strollur R package. Our
goal was to develop a framework to store amplicon sequence data that
could be integrated within other packages without the end user
necessarily knowing they were using strollur. This is meant to be
analogous to how many users of the tidyverse use the tibble package
without knowing they are using the package. By analogy, strollur strives
to provide a consistent interface to amplicon sequence data that also
integrates with analysis and visualization tools from the tidyverse and
other tools used by microbiome researchers.

strollur's data structures allows users to keep track of aligned and
unaligned sequence data, sequencing quality data, abundance data,
classification data, phylogenetic trees, operational taxonomic unit
(OTU) and amplicon sequence variant (ASV) assignments, and metadata.
Furthermore, the data structures can accommodate information describing
the commands that generated the data products and information about the
databases that were used. The package includes user-facing functions for
getting, assigning, and summarizing various forms of data. The package
also provides functions to read data from or write data to mothur
{[}\textbf{\emph{ref}}{]}, qiime2 {[}\textbf{\emph{ref}}{]}, dada2
{[}\textbf{\emph{ref}}{]}, phyloseq {[}\textbf{\emph{ref}}{]}, and biom
formats {[}\textbf{\emph{ref}}{]}. strollur can be used on any operating
system with no additional dependencies.

strollur also includes functions designed for developers to integrate
strollur objects into their own R packages. These developer-facing
functions provide additional access to the back end data. The ability to
import and export data as a strollur object has already been integrated
into the clustur, phylotypr, and rchime packages. The package leverages
the R6 object oriented programming framework and the Rcpp package to
accelerate data processing. By using these features of R, strollur
minimize the need to create deep copies of large data objects.

\textbf{Data availability.} strollur can accessed through the
Comprehensive R Archive Network (CRAN; DOI:
10.32614/CRAN.package.strollur) and developmental versions are available
through the project's GitHub website
(https://github.com/mothur2/strollur). The package's GitHub repository
includes the package's source code, the website's source code, automated
tests written using the testthat R package, and an issue tracker where
users can post questions, bug reports, and suggestions for future
features. A pkgdown version of the documentation including multiple
articles describing its use is hosted at (https://mothur.org/strollur).
The strollur package is available under the MIT License.

\subsection{Acknowledgements}\label{acknowledgements}

This project was supported, in part, by \ldots{}

\newpage

\subsection{References}\label{references}

\phantomsection\label{refs}




\end{document}
